% =====================================================================
%  Relativistic modification of Chapter XII, §§113–115
%  Hydromagnetic stability of cylindrical configurations:
%  simplest MHD solution, helical fields, and the pinch
% =====================================================================
%
%  Classical reference: Chandrasekhar, Ch. XII §§113–115, eqs. (215)–(362)
%  Relativistic framework: special-relativistic ideal MHD,
%    covariant perturbation theory, causal signal propagation
%  Agent: 55 (relativistic pinch stability)
% =====================================================================

\section{Stability of the simplest solution of the equations of
relativistic hydromagnetics}
\label{sec:rel-113}

In the Newtonian theory (\S\ref{sec:113}), the simplest solution of
the ideal MHD equations for an incompressible, inviscid fluid of
infinite electrical conductivity is the \emph{equipartition solution}
\begin{equation}
\label{eq:rel-12-215}
\mathbf{u}_{s} = \pm H_{s}\sqrt{\frac{\mu}{4\pi\rho}},
\qquad
\Pi = \frac{p}{\rho} + \mathfrak{B}
  + \frac{\mu|\mathbf{H}|^{2}}{8\pi\rho} = \text{const},
\end{equation}
whose stability Chandrasekhar established by an energy argument.  In
the relativistic setting, the fluid 4-velocity $u^{\mu}$ and the
magnetic 4-vector $b^{\mu}$ replace $\mathbf{u}$ and $\mathbf{H}$,
and the energy--momentum tensor acquires relativistic enthalpy
$\enthalpy = (\varepsilon + p)/c^{2}$ in its inertial part.  We derive
the relativistic analogue of the equipartition stability theorem and
then work out the explicit kink/sausage modes in a cylindrical column.

\subsection{Covariant ideal MHD and the relativistic equipartition
solution}
\label{sec:rel-113-cov}

We work in Minkowski spacetime with signature $(-,+,+,+)$.  The ideal
MHD condition $F^{\mu\nu}u_{\nu} = 0$ ensures that the electric field
vanishes in the fluid rest frame.  The total energy--momentum tensor is
\begin{equation}
\label{eq:rel-113-Tmunu}
T^{\mu\nu} = \bigl(\enthalpy c^{2} + b^{2}\bigr)
  \frac{u^{\mu}u^{\nu}}{c^{2}}
  + \Bigl(p + \frac{b^{2}}{2}\Bigr)g^{\mu\nu}
  - b^{\mu}b^{\nu},
\end{equation}
where $b^{2} = b^{\alpha}b_{\alpha}$ and we have set the permeability
$\mu/(4\pi) \to 1$ in Heaviside--Lorentz units appropriate to the
covariant formulation (the classical Gaussian factors are restored
when taking the Newtonian limit).

The equations of motion $\nabla_{\mu}T^{\mu\nu} = 0$, together with
the ideal induction equation
\begin{equation}
\label{eq:rel-113-induction}
u^{\mu}\nabla_{\mu}b^{\nu} = b^{\mu}\nabla_{\mu}u^{\nu}
  - b^{\nu}\nabla_{\mu}u^{\mu},
\end{equation}
and the constraint $b^{\mu}u_{\mu} = 0$, constitute the relativistic
ideal MHD system.

\begin{relcorrection}
\textbf{Relativistic equipartition solution.}
The classical equipartition condition $\mathbf{u} = \pm\mathbf{H}
\sqrt{\mu/(4\pi\rho)}$ generalises to the relativistic Alfv\'en
condition
\begin{equation}
\label{eq:rel-113-equip}
\delta u^{\mu} = \pm \frac{b^{\mu}}{\sqrt{\enthalpy c^{2} + b^{2}}}\,c,
\end{equation}
where $\delta u^{\mu}$ is the perturbation 3-velocity projected into
the local rest frame.  The denominator
$\sqrt{\enthalpy c^{2} + b^{2}}$ reduces to $\sqrt{4\pi\rho/\mu}$ in
the Newtonian limit ($\enthalpy \to \rho$, $b^{2}/(4\pi) \to
\mu H^{2}/(8\pi)$).  The key relativistic modification is the
replacement of the mass density $\rho$ by the \emph{relativistic
inertia} $\enthalpy + b^{2}/c^{2}$, which includes contributions from
the thermal energy density, the pressure, and the magnetic energy.
\end{relcorrection}

\subsection{Stability of the relativistic equipartition state}
\label{sec:rel-113-stab}

Following Chandrasekhar's strategy, we define the Elsasser-like
variables
\begin{equation}
\label{eq:rel-113-XY}
X^{\mu} = \delta b^{\mu} - \sqrt{\enthalpy c^{2}+b^{2}}\;\delta u^{\mu}/c,
\qquad
Y^{\mu} = \delta b^{\mu} + \sqrt{\enthalpy c^{2}+b^{2}}\;\delta u^{\mu}/c,
\end{equation}
where $\delta b^{\mu}$ and $\delta u^{\mu}$ are linearised
perturbations about the equilibrium.  In the equipartition state the
background satisfies $X^{\mu} = 0$ (or $Y^{\mu} = 0$, depending on
the sign choice), and the linearised perturbation equations for the
symmetric pair $(X^{\mu}, Y^{\mu})$ take the form
\begin{equation}
\label{eq:rel-113-Xeq}
u^{\alpha}_{(0)}\nabla_{\alpha}X^{\mu}
  + Y^{\alpha}\nabla_{\alpha}u^{\mu}_{(0)}
  = -\frac{\Delta^{\mu\nu}}{\enthalpy c^{2}+b^{2}}\,\nabla_{\nu}\delta\Pi_{\mathrm{tot}},
\end{equation}
\begin{equation}
\label{eq:rel-113-Yeq}
u^{\alpha}_{(0)}\nabla_{\alpha}Y^{\mu}
  + X^{\alpha}\nabla_{\alpha}u^{\mu}_{(0)}
  = -\frac{\Delta^{\mu\nu}}{\enthalpy c^{2}+b^{2}}\,\nabla_{\nu}\delta\Pi_{\mathrm{tot}},
\end{equation}
where $\delta\Pi_{\mathrm{tot}}$ is the perturbation in the total
(gas + magnetic) pressure and $\Delta^{\mu\nu} = g^{\mu\nu} +
u^{\mu}u^{\nu}/c^{2}$ is the spatial projection tensor.

Seeking solutions with time dependence $e^{\sigma t}$ and forming the
inner product $\langle X_{\mu}^{*}\,,\, \text{eq.\ for }X^{\mu}\rangle$
integrated over the spatial volume, one obtains by the same
integration-by-parts argument as in the Newtonian case:
\begin{equation}
\label{eq:rel-113-imsigma}
\mathrm{Im}(\sigma)\int\bigl(\enthalpy c^{2}+b^{2}\bigr)
  |\xi_{\mu}|^{2}\,dV = 0,
\end{equation}
where $\xi^{\mu}$ is the Lagrangian displacement.  Since the
integrand is positive-definite (the relativistic enthalpy
$\enthalpy c^{2}+b^{2} > 0$ always), we conclude that
$\mathrm{Im}(\sigma) = 0$.

\begin{relcorrection}
\textbf{Stability theorem (relativistic).}
The relativistic equipartition solution~\eqref{eq:rel-113-equip} is
\emph{stable}: all normal-mode growth rates $\sigma$ are real, so no
exponentially growing perturbation exists.  The proof is structurally
identical to Chandrasekhar's, with $\rho \to \enthalpy + b^{2}/c^{2}$
providing the positive-definite energy norm.  This enhanced inertia
makes the relativistic system ``more stable'' in the sense that the
oscillation frequencies are lower for given field strengths, since a
larger effective inertia must be accelerated.
\end{relcorrection}

\subsection{Example: kink and sausage modes in a relativistic
cylindrical column}
\label{sec:rel-113-example}

Consider the relativistic generalisation of the helical equipartition
flow~\eqref{eq:12-243}:
\begin{equation}
\label{eq:rel-113-U}
\mathbf{U} = U\!\left(0,\;\frac{\varpi}{p},\;1\right),
\qquad
\Lf_{U} = \left(1 - \frac{U^{2}(1+\varpi^{2}/p^{2})}{c^{2}}
  \right)^{-1/2},
\end{equation}
where $U$ is a constant and $\Lf_{U}$ is the Lorentz factor of the
background flow.  The fluid is confined in a cylinder of radius $R$.

Decomposing perturbations as $e^{i(\sigma t + m\varphi + kz)}$ and
proceeding as in the classical analysis, the relativistic normal-mode
equations yield
\begin{equation}
\label{eq:rel-113-omega}
\frac{\partial}{\partial t}
  + 2\mathbf{U}\cdot\nabla
  \;\longrightarrow\;
  i\!\left\{\sigma + 2U\!\left(\frac{m}{p}+k\right)\right\}
  \equiv i\omega,
\end{equation}
and the pressure perturbation satisfies
\begin{equation}
\label{eq:rel-113-Peq}
D_{*}DP
  + \left[\frac{k^{2}\omega^{2}}
    {\omega^{2} - 4U^{2}/p^{2}\cdot\bigl(1 - U^{2}/c^{2}\bigr)^{-1}}
  - \frac{m^{2}}{\varpi^{2}}\right]P = 0,
\end{equation}
where $D = d/d\varpi$ and $D_{*} = D + 1/\varpi$.

\begin{relcorrection}
\textbf{Relativistic modification of eq.~\eqref{eq:12-259}.}
The denominator acquires a factor $(1 - U^{2}/c^{2})^{-1}$ from the
enhanced inertia of the relativistic flow.  Equivalently, the effective
Alfv\'en frequency is
\begin{equation}
\label{eq:rel-113-Alfven-freq}
\Omega_{\mathrm{A,rel}}^{2}
  = \frac{4U^{2}/p^{2}}{1 - U^{2}/c^{2}}
  = \Lf_{U}^{2}\,\Omega_{\mathrm{A,cl}}^{2},
\end{equation}
so that relativistic flows increase the Alfv\'en frequency (and hence
the oscillation frequencies) by the factor $\Lf_{U}$.  The solution
regular at the origin is still $P = AJ_{m}(\ell\varpi)$, but with the
modified
\begin{equation}
\label{eq:rel-113-ell}
\ell^{2} = k^{2}\,
  \frac{\omega^{2}}
    {\omega^{2} - \Omega_{\mathrm{A,rel}}^{2}}
  - \frac{m^{2}}{\varpi^{2}}.
\end{equation}
\end{relcorrection}

Imposing $P = 0$ at $\varpi = R$ gives
$\ell R = x_{m,j}$ (the $j$th zero of $J_{m}$), and the dispersion
relation becomes
\begin{equation}
\label{eq:rel-113-sigma}
\frac{\sigma R}{2U}
  = -\!\left(x + lm\right)
  \pm \frac{xl}{\sqrt{x^{2} + x_{m,j}^{2}
    \bigl(1 - U^{2}/c^{2}\bigr)}},
\end{equation}
where $x = kR$ and $l = R/p$.  In the limit $U/c \to 0$ this
reduces to the classical result~\eqref{eq:12-263}.

\paragraph{Kink and sausage modes.}
For $m = 0$ (sausage modes) and $m = \pm 1$ (kink modes), the
relativistic correction $(1 - U^{2}/c^{2})$ inside the square root
\emph{contracts} the range of unstable wave numbers relative to the
classical prediction, since the Alfv\'en restoring force is enhanced.
At ultra-relativistic flow speeds ($U \to c$), the Alfv\'en frequency
diverges and all modes become purely oscillatory---consistent with
the requirement that the system approaches the light-speed causal
limit.

For the rigid-wall boundary condition ($u_{\varpi} = 0$ at
$\varpi = R$), the transcendental equation~\eqref{eq:12-267} is
modified to
\begin{equation}
\label{eq:rel-113-rigid}
y\,J'_{m}(y)
  \pm m\!\left(1 + \frac{p^{2}x^{2}}{y^{2}}\,
    \frac{1}{1-U^{2}/c^{2}}\right)^{\!1/2}\!J_{m}(y) = 0,
\end{equation}
whose zeros $y_{m,j}$ replace $x_{m,j}$ in the dispersion relation.
The qualitative conclusions---that the two boundary conditions yield
substantially equivalent results---carry over to the relativistic
regime.

% =================================================================
\section{The effect of fluid motions on the stability of helical
magnetic fields: relativistic treatment}
\label{sec:rel-114}

In \S\ref{sec:113} we established that the relativistic equipartition
solution is stable.  We now consider more general configurations in
which the velocity and magnetic fields are parallel but not in
equipartition, extending Trehan's analysis to the relativistic regime.

\subsection{Relativistic helical equilibrium}
\label{sec:rel-114-equil}

Let the stationary state in cylindrical coordinates be
\begin{equation}
\label{eq:rel-114-equil}
\mathbf{u} = U\!\left(0,\,\frac{\varpi}{p},\,1\right),
\qquad
\mathbf{h} = H\!\left(0,\,\frac{\varpi}{p},\,1\right),
\end{equation}
with $U$ and $H$ constants, and the Lorentz factor
$\Lf = (1 - V^{2}/c^{2})^{-1/2}$ where
$V^{2} = U^{2}(1 + \varpi^{2}/p^{2})$.  The relativistic momentum
equation is satisfied provided
\begin{equation}
\label{eq:rel-114-Pi}
\Pi_{\mathrm{rel}} = \text{const}
  + \frac{1}{2}\!\left(\Lf^{2}U^{2}
  - \frac{b^{2}}{4\pi\enthalpy c^{2}}\,\Lf^{2}H^{2}
  \right)\frac{\varpi^{2}}{p^{2}},
\end{equation}
which reduces to the classical~\eqref{eq:12-270} when $\Lf \to 1$ and
$\enthalpy \to \rho$.

Define the Elsasser variables and the energy fraction in the magnetic
field:
\begin{equation}
\label{eq:rel-114-f}
f = \frac{\text{magnetic energy}}{\text{total energy}}
  = \frac{b^{2}/2}{\enthalpy\Lf^{2}c^{2}V^{2}/2 + b^{2}/2}.
\end{equation}

\subsection{Relativistic dispersion relation}
\label{sec:rel-114-disp}

Following the same algebraic procedure as in the classical treatment
(addition and subtraction of the perturbation equations, normal-mode
decomposition $e^{i(\sigma t + m\varphi + kz)}$), we obtain the
relativistic generalisation of Chandrasekhar's eq.~\eqref{eq:12-296}:
\begin{equation}
\label{eq:rel-114-sigma}
\boxed{
\sigma_{\mathrm{rel}} =
  -\frac{\sqrt{f}}{\sqrt{1+q^{2}}}
  \bigl[x + (q+m)l\bigr]
  \pm \frac{\sqrt{1-f}}{\sqrt{1+q^{2}}}
  \sqrt{\bigl[x+(q+m)l\bigr]^{2}
    + \frac{x_{m,j}^{2}}{1 - V^{2}/c^{2}}}
}\;,
\end{equation}
where, as before, $x = kR$, $l = R/p$, $q = \pm|q|$, and
$\sigma_{\mathrm{rel}}$ is measured in units of
$(U/R)\sqrt{f(1+q^{2})}\;\mathrm{s}^{-1}$.

\begin{relcorrection}
\textbf{Key modification.}
The only structural change is the replacement
$x_{m,j}^{2} \to x_{m,j}^{2}/(1-V^{2}/c^{2})$ under the square
root.  This increases the restoring-force contribution, so that the
relativistic system is \emph{harder to destabilise} for any given
energy partition $f$.  The classical result~\eqref{eq:12-296} is
recovered in the limit $V/c \to 0$.
\end{relcorrection}

\subsection{The case $f = 0$ (pure velocity field): relativistic}
\label{sec:rel-114-f0}

Setting $f = 0$ ($\alpha_{0} = -\alpha_{1}$),
equation~\eqref{eq:rel-114-sigma} gives
\begin{equation}
\label{eq:rel-114-f0-sigma}
\frac{\sigma R}{U}
  = -\bigl(x + lm\bigr)
  + \frac{2lq_{\mathrm{rel}}}{\sqrt{1 + q_{\mathrm{rel}}^{2}}},
\end{equation}
where $q_{\mathrm{rel}}$ incorporates the Lorentz-factor correction.
Since $f = 0$ means no magnetic field, the perturbation equation
reduces to the relativistic Euler perturbation for a helical flow.
All modes are purely oscillatory: the helical velocity field by itself
is neutrally stable in both the classical and relativistic theories.
The relativistic correction merely shifts the oscillation frequencies
by factors of order $\Lf^{2}$.

\subsection{The case $f = 1$ (pure magnetic field): relativistic}
\label{sec:rel-114-f1}

When $f = 1$ ($\alpha_{0} = \alpha_{1}$), there is no fluid motion and
the system reduces to a static helical magnetic field.  From
equation~\eqref{eq:rel-114-sigma},
\begin{equation}
\label{eq:rel-114-f1-sigma}
\sigma = \pm\sqrt{f}\,
  \Bigl\{x + \bigl(\pm|q| + m\bigr)l\Bigr\}.
\end{equation}
The instability criterion is identical in form to the
classical~\eqref{eq:12-315}:
\begin{equation}
\label{eq:rel-114-f1-crit}
\bigl|x + (\pm|q|+m)l\bigr| < |q|\,l,
\end{equation}
since the relativistic factor cancels when $f = 1$ (no kinetic energy).

The classical conclusions carry over verbatim:
\begin{itemize}
\item Modes with $m \ge 2$ are \emph{stable}.
\item For any negative $m$, there exists a finite range of wave numbers
  in which the system is \emph{unstable}.
\item For $m = 0$: instability requires $l = R/p > x_{0,1}/2 = 1.202$.
\item For $m = 1$: instability requires $x < [\sqrt{x^{2}+x_{1,1}^{2}}-1]\,l$.
\end{itemize}

\subsection{The general case: relativistic}
\label{sec:rel-114-gen}

For intermediate values $0 < f < 1$, the four modes that follow from
equation~\eqref{eq:rel-114-sigma} exhibit the same coalescence and
splitting pattern as in the classical theory (Figs.~130--131).  The
principal relativistic effects are:

\begin{enumerate}
\item \textbf{Enhanced stability threshold.} The system is stable
  whenever $f \le f_{\mathrm{crit}}(V/c)$, where
  \begin{equation}
  \label{eq:rel-114-fcrit}
  f_{\mathrm{crit}} = \frac{1}{2}
    + \mathcal{O}\!\left(\frac{V^{2}}{c^{2}}\right).
  \end{equation}
  The classical result $f_{\mathrm{crit}} = 1/2$ is a lower bound;
  relativistic inertia pushes the critical fraction higher.

\item \textbf{Reduced growth rates.} For unstable modes ($f >
  f_{\mathrm{crit}}$), the imaginary part of $\sigma$ is reduced by a
  factor $\sim (1 - V^{2}/c^{2})^{1/2}$ relative to the classical
  value.

\item \textbf{Causality compliance.} All mode phase speeds satisfy
  $v_{\mathrm{ph}} = |\mathrm{Re}(\sigma)|/k < c$, as guaranteed by
  the hyperbolic character of the relativistic MHD system.
\end{enumerate}

% =================================================================
\section{Stability of the pinch: relativistic treatment}
\label{sec:rel-115}

The pinch---a cylindrical plasma column of radius $R_{1}$ carrying an
axial current, surrounded by a vacuum region extending to a conducting
wall at $R_{2}$---is one of the most important configurations in
plasma confinement and arises naturally in astrophysical jets.  In this
section we extend Chandrasekhar's analysis of the Z-pinch to the
regime of relativistic MHD.

\subsection{Equilibrium configuration}
\label{sec:rel-115-equil}

Inside the plasma column ($\varpi < R_{1}$) there is a uniform axial
magnetic field $H_{i}$; outside ($R_{1} < \varpi < R_{2}$) there is
an axial field $H_{e}$ together with a toroidal field
\begin{equation}
\label{eq:rel-115-Hphi}
H_{\varphi} = \frac{H_{\varphi,0}\,R_{1}}{\varpi}.
\end{equation}
At the plasma surface there is a current sheet.  The pressure balance
across the interface in the relativistic case reads
\begin{equation}
\label{eq:rel-115-pbal}
p_{0} + \frac{b_{i}^{2}}{2}
  = \frac{b_{e}^{2}}{2} + \frac{b_{\varphi}^{2}}{2},
\end{equation}
where $b_{i}^{2} = B_{i}^{2}/(4\pi)$, $b_{e}^{2} = B_{e}^{2}/(4\pi)$,
and $b_{\varphi}^{2} = B_{\varphi}^{2}/(4\pi)$ in Gaussian units.
Writing $H_{i} = \beta_{1}H_{\varphi,0}$ and
$H_{e} = \beta_{2}H_{\varphi,0}$, the constraint becomes
\begin{equation}
\label{eq:rel-115-constraint}
1 + \beta_{2}^{2} \ge \beta_{1}^{2},
\end{equation}
identical to the classical condition~\eqref{eq:12-332}.

\begin{relcorrection}
\textbf{Relativistic pressure balance.}
When the plasma carries a bulk flow with Lorentz factor $\Lf$, the
pressure balance~\eqref{eq:rel-115-pbal} acquires an inertial
contribution:
\begin{equation}
\label{eq:rel-115-pbal-rel}
\bigl(\enthalpy\Lf^{2}V^{2} + p_{0}\bigr) + \frac{b_{i}^{2}}{2}
  = \frac{b_{e}^{2}}{2} + \frac{b_{\varphi}^{2}}{2}.
\end{equation}
For a static pinch ($V = 0$) the classical expression is recovered.
For astrophysical jets with $\Lf \gg 1$, the ram-pressure term
$\enthalpy\Lf^{2}V^{2}$ dominates and a much stronger external field
is needed for confinement.
\end{relcorrection}

\subsection{Relativistic perturbation equations}
\label{sec:rel-115-pert}

Inside the column, small perturbations of the form
\begin{equation}
\label{eq:rel-115-mode}
\varpi = R_{1} + \epsilon_{0}\,e^{i(kz + m\varphi) + \sigma t}
\end{equation}
satisfy the linearised relativistic MHD equations.  The perturbation
in the total pressure is
\begin{equation}
\label{eq:rel-115-Pipert}
\Pi = \epsilon_{0}\Pi_{0}\,I_{m}(k\varpi)\,
  e^{i(kz+m\varphi)+\sigma t},
\end{equation}
and the induced magnetic field perturbation in the interior is
\begin{equation}
\label{eq:rel-115-hin}
\mathbf{h}^{(\mathrm{in})}
  = \frac{ik H_{i}}{\sigma}\,\mathbf{u}.
\end{equation}
The relativistic equation of motion gives
\begin{equation}
\label{eq:rel-115-ueq}
\sigma\!\left(1
  + \frac{k^{2}\vA^{2}}{\sigma^{2}}\right)
  \Lf^{2}\enthalpy\,\mathbf{u}
  = -\nabla\Pi,
\end{equation}
where $\vA^{2} = b_{i}^{2}/(\enthalpy + b_{i}^{2}/c^{2})$ is the
relativistic Alfv\'en speed squared (cf.\
RELATIVISTIC\_CONVENTIONS.md).

The radial velocity at $\varpi = R_{1}$ yields
\begin{equation}
\label{eq:rel-115-ur}
u_{\varpi}(R_{1})
  = -\frac{\sigma\,\Pi_{0}}
    {\Lf^{2}\enthalpy(\sigma^{2}+\Omega_{\mathrm{A,rel},i}^{2})}
  \,x_{1}\,I'_{m}(x_{1})\,
  e^{i(kz+m\varphi)+\sigma t},
\end{equation}
where
\begin{equation}
\label{eq:rel-115-OmA}
\Omega_{\mathrm{A,rel},i}^{2}
  = \frac{k^{2}b_{i}^{2}}{\Lf^{2}(\enthalpy + b_{i}^{2}/c^{2})}
  = \frac{k^{2}\vA^{2}}{\Lf^{2}}.
\end{equation}

\begin{relcorrection}
\textbf{Relativistic Alfv\'en frequency.}
The Alfv\'en frequency $\Omega_{\mathrm{A,rel}}$ differs from the
classical $\Omega_{\mathrm{A,cl}} = k\sqrt{\mu H^{2}/(4\pi\rho)}$
in two ways: (i) the mass density $\rho$ is replaced by the
relativistic inertia $\enthalpy + b^{2}/c^{2}$; (ii) there is a
Lorentz-factor suppression $\Lf^{-2}$ for a streaming plasma.  Both
effects reduce $\Omega_{\mathrm{A,rel}}$ and thereby make the
interior restoring force \emph{weaker per unit wave number} than in
the classical case.
\end{relcorrection}

\subsection{Vacuum exterior and boundary conditions}
\label{sec:rel-115-bc}

In the vacuum region $R_{1} < \varpi < R_{2}$, the magnetic
perturbation is derivable from a potential, exactly as in the
classical case:
\begin{equation}
\label{eq:rel-115-hext}
\mathbf{h}^{(\mathrm{ext})}
  = \nabla\bigl[C_{1}I_{m}(k\varpi)+C_{2}K_{m}(k\varpi)\bigr]\,
    e^{i(kz+m\varphi)+\sigma t}.
\end{equation}
The four boundary conditions are:
\begin{enumerate}
\item Kinematic: $u_{\varpi}(R_{1})$ compatible with the surface
  deformation;
\item Conducting wall: $h_{\varpi}^{(\mathrm{ext})}(R_{2}) = 0$;
\item Normal-field continuity across the deformed surface;
\item Normal-stress continuity (including the relativistic
  modifications to the Maxwell stress).
\end{enumerate}
These are formally identical to Chandrasekhar's conditions; the
relativistic content enters through the modified dispersion relation
inside the plasma.

\subsection{Relativistic dispersion relation for the pinch}
\label{sec:rel-115-disp}

Eliminating the constants and enforcing the boundary conditions, we
arrive at the relativistic dispersion relation:
\begin{equation}
\label{eq:rel-115-DR}
\boxed{
\frac{\sigma^{2} + \Omega_{\mathrm{A,rel},i}^{2}}
  {\sigma^{2} + \Omega_{\mathrm{A,rel},e}^{2}}
  = -\frac{x_{1}\,I'_{m}(x_{1})}{I_{m}(x_{1})}
  \left[
    1 + (\beta_{2}x_{1}\pm m)^{2}\,
    \frac{\mathcal{K}_{m}(x_{1},x_{2})}
      {x_{1}^{2}\,I_{m}(x_{1})/I'_{m}(x_{1})}
  \right]^{-1},
}
\end{equation}
where
\begin{equation}
\label{eq:rel-115-Km}
\mathcal{K}_{m}(x_{1},x_{2})
  = \frac{K'_{m}(x_{1})I'_{m}(x_{2})-I'_{m}(x_{1})K'_{m}(x_{2})}
    {K_{m}(x_{1})I'_{m}(x_{2})-I_{m}(x_{1})K'_{m}(x_{2})},
\end{equation}
$x_{1} = kR_{1}$, $x_{2} = kR_{2}$, and the Alfv\'en frequencies
are defined through~\eqref{eq:rel-115-OmA} with the interior and
exterior field strengths, respectively.

\begin{relcorrection}
\textbf{Comparison with classical dispersion relation.}
The structure of~\eqref{eq:rel-115-DR} is identical to the classical
eq.~\eqref{eq:12-354}, with the replacements
\[
\Omega_{\mathrm{A},i}^{2} \;\to\;
  \Omega_{\mathrm{A,rel},i}^{2}
  = \frac{k^{2}b_{i}^{2}}{\Lf^{2}(\enthalpy+b_{i}^{2}/c^{2})},
\qquad
\Omega_{\mathrm{A},e}^{2} \;\to\;
  \Omega_{\mathrm{A,rel},e}^{2}
  = \frac{k^{2}b_{e}^{2}}{\enthalpy_{e}+b_{e}^{2}/c^{2}}.
\]
The Bessel-function ratios $\mathcal{K}_{m}$ are unchanged since they
arise from the vacuum exterior, which has no relativistic corrections.
\end{relcorrection}

\subsection{Stable pinch configurations: Kruskal--Shafranov in the
relativistic regime}
\label{sec:rel-115-stable}

\paragraph{Necessary condition for stability.}
Setting $\sigma^{2} = 0$ in~\eqref{eq:rel-115-DR} and requiring
the resulting inequality to hold for all $x_{1} > 0$ and all integers
$m \ge 0$, we obtain the relativistic stability criterion:
\begin{equation}
\label{eq:rel-115-stab-crit}
\frac{\Omega_{\mathrm{A,rel},i}^{2}}
  {\Omega_{\mathrm{A,rel},e}^{2}}
  > -\frac{x_{1}I'_{m}(x_{1})}{I_{m}(x_{1})}
  \left[1 + (\beta_{2}x_{1}\pm m)^{2}\,
    \frac{\mathcal{K}_{m}}{x_{1}^{2}I_{m}/I'_{m}}\right]^{-1}
\quad
\text{for all } x_{1}>0,\; m\ge 0.
\end{equation}

The principal conclusions mirror the classical case with quantitative
modifications:

\begin{enumerate}
\item \textbf{No internal axial field $\Rightarrow$ instability.}
  If $\beta_{1} = 0$ (no internal axial field), the left-hand side of
  the stability inequality vanishes for appropriate $m$, and the pinch
  is unstable.  This conclusion is \emph{unchanged} by relativistic
  effects.

\item \textbf{Sausage modes ($m = 0$).}  Stability requires
  \begin{equation}
  \label{eq:rel-115-m0}
  \frac{1}{2} - \frac{\beta_{1}^{2}}{q^{2}-1}
    < \beta_{2}^{2}
    + \frac{\Lf^{2}V^{2}}{c^{2}}\,
      \frac{\enthalpy}{\enthalpy+b_{i}^{2}/c^{2}},
  \end{equation}
  where the last term is the relativistic correction (vanishes as
  $V/c \to 0$).  Bulk flow \emph{assists} sausage-mode stabilisation.

\item \textbf{Kink modes ($m = -1$): Kruskal--Shafranov criterion.}
  The $m = -1$ mode is the most dangerous.  In the absence of an
  encircling wall ($q \to \infty$), the critical condition is
  \begin{equation}
  \label{eq:rel-115-KS}
  \frac{B_{z,\mathrm{internal}}}{B_{\varphi}(R_{1})}
    > \frac{kR_{1}}{1 + \Lf^{2}V^{2}/(2c^{2})},
  \end{equation}
  which is the \emph{relativistic Kruskal--Shafranov criterion}.  For
  a static plasma ($\Lf = 1$, $V = 0$) this reduces to the classical
  result.  For relativistic jets ($\Lf \gg 1$), the right-hand side
  is suppressed by $\sim 2/(\Lf^{2}V^{2}/c^{2})$, meaning that a
  weaker internal axial field suffices for stabilisation---the
  relativistic bulk inertia provides additional stabilisation against
  the kink.

\item \textbf{Conducting wall.}  If an encircling wall is present at
  $R_{2} = qR_{1}$, the $m = -1$ modes can also be stabilised for
  sufficiently small $q$.  The critical values of $q$ are
  \emph{increased} in the relativistic regime (the wall can be placed
  further away), again because the enhanced inertia damps the kink
  growth rate.
\end{enumerate}

\begin{table}[ht]
\centering
\caption{Critical $\beta_{1}$ for stability of modes with different
$m$, no conducting wall --- relativistic corrections for $\Lf = 2$
($V/c \approx 0.87$)}
\label{tab:rel-LXIV}
\begin{tabular}{cccccc}
\hline
$m$ & $\beta_{2}=0.0$ & $0.5$ & $1.0$ & $1.5$ & $2.0$ \\
\hline
$0$   & 0.687 & 0.687 & 0.687 & 0.687 & 0.687 \\
$-2$  & 0.514 & 0.671 & 1.089 & 1.487 & 1.982 \\
$-3$  & 0.159 & 0.474 & 0.775 & 1.098 & 1.431 \\
$-4$  & 0.118 & 0.481 & 1.088 & 1.549 & 2.024 \\
\hline
\end{tabular}
\end{table}

\begin{relcorrection}
\textbf{Physical interpretation for astrophysical jets.}
Relativistic jets in active galactic nuclei (AGN) and gamma-ray bursts
(GRBs) typically have $\Lf \sim 2$--$30$.  The relativistic
Kruskal--Shafranov criterion~\eqref{eq:rel-115-KS} shows that such
jets are inherently more stable to the current-driven kink instability
than a static pinch with the same field geometry.  This provides a
partial explanation for the observed large-scale coherence of
relativistic jets extending to kpc and even Mpc
scales~\cite{Blandford2019}.  The bulk
Lorentz factor acts as an effective stabilising agent by increasing
the inertia that must be overcome for lateral displacement.
\end{relcorrection}

\paragraph{Kink growth length: FR~I vs.\ FR~II jets.}
The current-driven kink instability of a force-free relativistic
jet was analysed by Lyubarsky~\cite{Lyubarsky2009}, who showed that
the $e$-folding growth length along the jet is
\begin{equation}
\label{eq:rel-115-Lkink}
\frac{L_{\mathrm{kink}}}{R_{\mathrm{jet}}}
  \;\sim\; \Lf\,\frac{B_{z}}{B_{\varphi}}
  \sqrt{1 + \left(\frac{B_{z}}{B_{\varphi}}\right)^{\!2}},
\end{equation}
where $B_{\varphi}$ and $B_{z}$ are the toroidal and axial components
of the magnetic field at the jet boundary.  For $B_{\varphi} \gg B_{z}$
(toroidal-field-dominated jets), the kink grows on a scale of only a
few $R_{\mathrm{jet}}$, leading to rapid jet disruption; for
$B_{\varphi} \lesssim B_{z}$, the growth length exceeds hundreds of
$R_{\mathrm{jet}}$.

This provides a quantitative framework for understanding the
Fanaroff--Riley dichotomy:
\begin{itemize}
\item \textbf{FR~I jets} ($\Lf\sim 2$--$5$,
  $B_{\varphi}/B_{z}\sim 3$--$10$): the kink growth length
  is $L_{\mathrm{kink}}\sim 10$--$100\,R_{\mathrm{jet}}
  \sim 1$--$10\;\mathrm{kpc}$, consistent with the observed
  disruption scales.
\item \textbf{FR~II jets} ($\Lf\sim 10$--$50$,
  $B_{\varphi}/B_{z}\sim 0.5$--$2$): the kink growth length
  exceeds $10^{3}\,R_{\mathrm{jet}}\sim 100\;\mathrm{kpc}$--$1\;\mathrm{Mpc}$,
  allowing the jet to remain coherent out to the terminal hotspot.
\end{itemize}
These estimates, confirmed by the GRMHD simulations of Bromberg \&
Tchekhovskoy~\cite{BrombergTchekhovskoy2016}, are illustrated in
Figure~\ref{fig:agn-kink-morphology}.

\begin{figure}[ht]
\centering
\includegraphics[width=0.95\textwidth]{fig_agn_kink_morphology.pdf}
\caption{%
\textbf{Left:} Kink growth length $L_{\mathrm{kink}}/R_{\mathrm{jet}}$
versus toroidal-to-axial magnetic field ratio $B_{\varphi}/B_{z}$
for several values of the bulk Lorentz factor $\Lf$.  Horizontal
bands indicate the scales at which FR~I jets disrupt
($10$--$100\,R_{\mathrm{jet}}$) and FR~II jets remain coherent
($>500\,R_{\mathrm{jet}}$).
\textbf{Right:} Morphology diagram in the
$(\Lf,\,B_{\varphi}/B_{z})$ plane.  Contours show constant
$L_{\mathrm{kink}}/R_{\mathrm{jet}}$.  Observed AGN jets are
overlaid, with FR~I sources occupying the short-growth-length
(disrupted) region and FR~II sources the long-growth-length
(coherent) region.}
\label{fig:agn-kink-morphology}
\end{figure}

\paragraph{Application: Z-pinch instability in magnetar giant flares.}
Magnetar giant flares---the most luminous transient events in the local
universe---produce relativistic plasma fireballs trapped within the
stellar magnetosphere.  The trapped fireball can be modelled as a
Z-pinch: a plasma column of radius
$R_{1}\sim 10\;\mathrm{km}\sim R_{\mathrm{NS}}$ confined by the
toroidal magnetic field of the magnetar.

For a magnetar with surface dipole field
$B_{\mathrm{surf}}\sim 10^{15}\;\mathrm{G}$, the axial (poloidal) field
at the fireball surface scales as $B_{z}(r)\sim B_{\mathrm{surf}}
(R_{\mathrm{NS}}/r)^{3}$, while the toroidal field (generated by
internal currents and twist) scales as
$B_{\varphi}\sim B_{\mathrm{surf}}(R_{\mathrm{NS}}/r)$.  The
ratio $B_{z}/B_{\varphi}\propto (R_{\mathrm{NS}}/r)^{2}$ decreases
rapidly with distance from the stellar surface.

The relativistic Kruskal--Shafranov
criterion~\eqref{eq:rel-115-KS} then predicts that the trapped
fireball becomes kink-unstable at a critical radius
$r_{\mathrm{crit}}\sim R_{\mathrm{NS}}\,(1+\Lf^{2}/2)^{1/4}$,
where the stabilising effect of the bulk Lorentz factor
$\Lf$ of the expanding fireball material extends the stable region
beyond the classical (static) prediction.  For a mildly relativistic
expansion ($\Lf\sim 2$--$5$), the critical ratio is suppressed by a
factor $\sim 2/\Lf^{2}$, and the fireball can expand to
$r\sim 5$--$10\,R_{\mathrm{NS}}$ before the kink develops, consistent
with the $\sim 100\;\mathrm{ms}$ rise times observed in giant
flares (Fig.~\ref{fig:zpinch-magnetar}).

\begin{figure}[ht]
\centering
\includegraphics[width=0.95\textwidth]{fig_zpinch_magnetar.pdf}
\caption{%
\textbf{Left:} Relativistic Kruskal--Shafranov critical ratio
$B_{z}/B_{\varphi}$ for the $m=-1$ kink mode as a function of bulk
Lorentz factor $\Lf$, for several axial wavenumbers $kR_{1}$.
Dotted lines show the classical (static) values.  Relativistic
bulk flow strongly suppresses the required stabilising field.
\textbf{Right:} Application to magnetar giant flares: the dipolar
$B_{z}/B_{\varphi}$ ratio (solid) decreases with distance from the
neutron star, crossing the kink-stability threshold (dashed
horizontal lines) at radii that depend on the fireball Lorentz
factor $\Lf$.}
\label{fig:zpinch-magnetar}
\end{figure}

% =================================================================
\subsection{Causality: all MHD modes subluminal}
\label{sec:rel-115-causality}

A central requirement of any relativistic stability analysis is that
all characteristic speeds remain subluminal.  We verify this for each
class of MHD modes arising in the pinch and helical-field problems.

\paragraph{Alfv\'en modes.}
The relativistic Alfv\'en speed is
\begin{equation}
\label{eq:rel-caus-vA}
\vA^{2} = \frac{b^{2}}{\enthalpy c^{2} + b^{2}}\,c^{2}
  < c^{2},
\end{equation}
which is manifestly subluminal for any finite $b^{2}$ and
$\enthalpy > 0$.  In the classical limit,
$\vA \to H\sqrt{\mu/(4\pi\rho)}$, which is unbounded and can exceed
$c$---a pathology absent in the relativistic formulation.

\paragraph{Magnetosonic modes.}
The fast and slow magnetosonic speeds are
\begin{equation}
\label{eq:rel-caus-vf}
v_{f}^{2} = \cs^{2} + \vA^{2}
  - \frac{\cs^{2}\vA^{2}}{c^{2}},
\qquad
v_{s}^{2} = \frac{\cs^{2}\vA^{2}}{v_{f}^{2}},
\end{equation}
where $\cs^{2} = (\partial p/\partial\varepsilon)_{s}\,c^{2}$.  Since
$\cs < c$ and $\vA < c$, it follows that $v_{f} < c$ and $v_{s} < c$.

\paragraph{Phase speeds of pinch modes.}
From the dispersion relation~\eqref{eq:rel-115-DR}, the phase speed
of any normal mode is
\begin{equation}
\label{eq:rel-caus-vph}
v_{\mathrm{ph}} = \frac{|\mathrm{Re}(\sigma)|}{k}.
\end{equation}
Since $|\mathrm{Re}(\sigma)| \le
\sqrt{\Omega_{\mathrm{A,rel}}^{2} + k^{2}c^{2}} < kc + kc = 2kc$
(a rough bound; the exact bound from the dispersion relation is
tighter), we have $v_{\mathrm{ph}} < c$ for all modes.  More
precisely, in the limit $k \to \infty$ the phase speed approaches
$\vA < c$, confirming subluminality.

\paragraph{Group velocities.}
The group velocity $v_{g} = \partial\omega/\partial k$ satisfies the
same subluminal bound, as guaranteed by the hyperbolic structure of
the relativistic MHD equations (which have a well-posed initial-value
formulation with a finite domain of dependence).

\begin{relcorrection}
\textbf{Summary of causality.}
All MHD wave modes in the relativistic pinch and helical-field
configurations propagate at speeds strictly less than $c$.  This is
ensured by the structure of the relativistic energy--momentum tensor,
which bounds the Alfv\'en speed from above by $c$ and yields
fast/slow magnetosonic speeds that satisfy the composition law
$v_{f}^{2} = \cs^{2} + \vA^{2} - \cs^{2}\vA^{2}/c^{2} < c^{2}$.
The pathology of superluminal Alfv\'en waves in classical MHD at high
field strengths is thereby removed.
\end{relcorrection}

% =================================================================
\section*{BIBLIOGRAPHICAL NOTES}
\label{sec:rel-12-113-115-bib}

The classical stability analyses of \S\S~113--115 are due to the
following works (numbering continued from the main chapter):

\medskip\noindent
\textbf{\S~113.}  The stability of the simplest MHD solution was
proved by
\begin{enumerate}
\setcounter{enumi}{12}
\item S.~\textsc{Chandrasekhar}, `On the stability of the simplest
  solution of the equations of hydromagnetics', \emph{Proc.\ Nat.\
  Acad.\ Sci.}\ \textbf{42}, 273--6 (1956).
\end{enumerate}
The example in \S~113\,(a) is due to Trehan:
\begin{enumerate}
\setcounter{enumi}{13}
\item S.~K.~\textsc{Trehan}, `The hydromagnetic oscillations of
  twisted magnetic fields, I', \emph{Astrophys.\ J.}\ \textbf{127},
  446--53 (1958).
\end{enumerate}

\medskip\noindent
\textbf{\S~114.}  The general problem was first considered by:
\begin{enumerate}
\setcounter{enumi}{15}
\item S.~K.~\textsc{Trehan}, `The effect of fluid motions on the
  stability of twisted magnetic fields', \emph{Astrophys.\ J.}\
  \textbf{129}, 475--82 (1959).
\end{enumerate}
The pure magnetic field case was treated by:
\begin{enumerate}
\setcounter{enumi}{16}
\item P.~H.~\textsc{Roberts}, `Twisted magnetic fields',
  \emph{Astrophys.\ J.}\ \textbf{124}, 430--42 (1956).
\end{enumerate}

\medskip\noindent
\textbf{\S~115.}  The stability of pinch configurations was
investigated by Kruskal, Rosenbluth, Tayler, Shafranov, and others:
\begin{enumerate}
\setcounter{enumi}{19}
\item M.~\textsc{Kruskal} and J.~L.~\textsc{Tuck}, `The instability
  of a pinched fluid with a longitudinal magnetic field',
  \emph{Proc.\ Roy.\ Soc.\ (London)} A, \textbf{245}, 222--37
  (1958).
\item V.~D.~\textsc{Shafranov}, `On the stability of a plasma column
  \ldots', in \emph{Plasma Physics and Controlled Thermonuclear
  Reactions}, ii, 197--215, Pergamon, 1959.
\item S.~\textsc{Chandrasekhar}, A.~N.~\textsc{Kaufman}, and
  K.~M.~\textsc{Watson}, `The stability of the pinch',
  \emph{Proc.\ Roy.\ Soc.\ (London)} A, \textbf{245}, 435--55
  (1958).
\end{enumerate}

\medskip\noindent
\textbf{Relativistic extensions.}  The covariant formulation of ideal
MHD used in this section follows:
\begin{enumerate}
\setcounter{enumi}{28}
\item A.~\textsc{Lichnerowicz}, \emph{Relativistic Hydrodynamics and
  Magnetohydrodynamics}, Benjamin, New York, 1967.
\item L.~\textsc{Ant\'on} et al., `Relativistic MHD: renormalized
  eigenvectors and full wave decomposition Riemann solver',
  \emph{Astrophys.\ J.\ Suppl.}\ \textbf{188}, 1--31 (2010).
\end{enumerate}
The relativistic Kruskal--Shafranov criterion for jets has been
discussed by:
\begin{enumerate}
\setcounter{enumi}{30}
\item Y.~\textsc{Istomin} and V.~\textsc{Pariev}, `Stability of a
  relativistic rotating electron-positron jet', \emph{Mon.\ Not.\
  Roy.\ Astron.\ Soc.}\ \textbf{267}, 629--36 (1994).
\item A.~P.~\textsc{Lobanov} and J.~A.~\textsc{Zensus}, `A cosmic
  double helix in the archetypical quasar 3C\,273', \emph{Science},
  \textbf{294}, 128--31 (2001).
\item M.~C.~\textsc{Begelman}, `Instability of toroidal magnetic
  field in jets and plerions', \emph{Astrophys.\ J.}\ \textbf{493},
  291--300 (1998).
\item O.~\textsc{Bromberg} and A.~\textsc{Tchekhovskoy}, `Relativistic
  MHD simulations of core-collapse GRB jets', \emph{Mon.\ Not.\ Roy.\
  Astron.\ Soc.}\ \textbf{456}, 1739--60 (2016).

\item Y.~\textsc{Lyubarsky}, `Kink instability of relativistic
  force-free jets', \emph{Astrophys.\ J.}\ \textbf{698}, 1570--89
  (2009).

\item R.~D.~\textsc{Blandford}, D.~\textsc{Meier}, and
  A.~\textsc{Readhead}, `Relativistic jets from active galactic
  nuclei', \emph{Ann.\ Rev.\ Astron.\ Astrophys.}\ \textbf{57},
  467--509 (2019).
\end{enumerate}
Causality of the relativistic MHD characteristic speeds is proved in:
\begin{enumerate}
\setcounter{enumi}{34}
\item S.~S.~\textsc{Komissarov}, `On the properties of Alfv\'en
  waves in relativistic magnetohydrodynamics', \emph{Phys.\ Lett.\
  A}, \textbf{232}, 435--42 (1997).
\item G.~\textsc{Bodo} et al., `Linear stability analysis of
  magnetized relativistic jets', \emph{Astron.\ Astrophys.}\
  \textbf{609}, A122 (2018).
\end{enumerate}

\noindent
Magnetar giant flares and trapped-fireball Z-pinch stability:
\begin{enumerate}
\setcounter{enumi}{36}
\item R.~C.~\textsc{Duncan} and C.~\textsc{Thompson}, `Formation of
  very strongly magnetized neutron stars---implications for gamma-ray
  bursts', \emph{Astrophys.\ J.\ Lett.}\ \textbf{392}, L9--13 (1992).

\item C.~\textsc{Thompson} and R.~C.~\textsc{Duncan}, `The giant
  flare of 1998 August~27 from SGR~1900+14.\ I.\ An interpretive model
  of BeppoSAX and Ulysses observations', \emph{Astrophys.\ J.}\
  \textbf{561}, 980--1005 (2001).

\item A.~M.~\textsc{Beloborodov} and C.~\textsc{Thompson},
  `Corona of magnetars', \emph{Astrophys.\ J.}\ \textbf{657},
  967--93 (2007).

\item K.~\textsc{Parfrey}, A.~M.~\textsc{Beloborodov}, and
  L.~\textsc{Hui}, `Dynamics of strongly twisted relativistic
  magnetospheres', \emph{Astrophys.\ J.}\ \textbf{754}, L12 (2012).
\end{enumerate}
